\documentclass{article}
\usepackage{amsmath,amssymb,amsthm}
\usepackage{algorithm,algorithmic}
\usepackage{hyperref}
\newtheorem{theorem}{Theorem}
\newtheorem{lemma}{Lemma}
\newtheorem{assumption}{Assumption}
\newtheorem{corollary}{Corollary}
\newcommand{\E}{\mathbb{E}}
\newcommand{\R}{\mathbb{R}}
\newcommand{\norm}[1]{\left\|#1\right\|}
\newcommand{\grad}{\nabla}
\begin{document}

\section*{Algorithm Name}
\textbf{SAGA (Stochastic Average Gradient Augmented).}

\section*{Mathematical Formulation}
Let the objective be a finite sum
\[
f(x)\;=\;\frac1n\sum_{i=1}^{n} f_i(x),\qquad f_i:\R^d\to\R .
\]
For each component $i$ we keep a \emph{memory point} $\phi_i^k\in\R^d$ at iteration $k$.
Define the stored gradient table
\[
g_i^k \;:=\; \grad f_i(\phi_i^k),\qquad 
\bar g^k \;:=\;\frac1n\sum_{i=1}^{n} g_i^k .
\]

Given a stepsize $\alpha>0$, the SAGA update at iteration $k$ proceeds as follows.
\begin{enumerate}
\item Sample an index $i_k$ uniformly from $\{1,\dots,n\}$.
\item Compute the \emph{variance‑reduced stochastic gradient}
\[
v_k \;:=\; \grad f_{i_k}(x_k)\;-\;\grad f_{i_k}(\phi_{i_k}^k)\;+\;\bar g^k .
\]
\item Update the iterate
\[
x_{k+1}\;=\;x_k\;-\;\alpha\, v_k . \tag{1}
\]
\item Update the table for the sampled component only:
\[
\phi_{i_k}^{k+1}\;=\;x_k,\qquad 
g_{i_k}^{k+1}\;=\;\grad f_{i_k}(x_k),
\]
and keep $\phi_j^{k+1}=\phi_j^{k},\;g_j^{k+1}=g_j^{k}$ for all $j\neq i_k$.
\end{enumerate}
The algorithm terminates after a predetermined number of epochs or when a
stopping criterion (e.g. $\norm{\grad f(x_k)}\le\varepsilon$) is met.

\section*{Pseudocode}
\begin{algorithm}[H]
\caption{SAGA}
\begin{algorithmic}[1]
\STATE \textbf{Input:} stepsize $\alpha>0$, number of iterations $K$.
\STATE Initialize $x_0\in\R^d$, choose $\phi_i^0$ (often $x_0$) for $i=1,\dots,n$.
\FOR{$k=0,\dots,K-1$}
    \STATE Sample $i_k\sim\text{Uniform}\{1,\dots,n\}$.
    \STATE $v_k\gets \grad f_{i_k}(x_k)-\grad f_{i_k}(\phi_{i_k}^k)+\frac1n\sum_{j=1}^{n}\grad f_j(\phi_j^k)$.
    \STATE $x_{k+1}\gets x_k-\alpha v_k$.
    \STATE $\phi_{i_k}^{k+1}\gets x_k$.
    \STATE \textbf{for} $j\neq i_k$: $\phi_j^{k+1}\gets\phi_j^{k}$.
\ENDFOR
\STATE \textbf{Output:} $x_K$ (or the average $\frac1K\sum_{k=0}^{K-1}x_k$).
\end{algorithmic}
\end{algorithm}

\section*{Dry Run Example}
Consider a scalar problem $f(w)=(w-3)^2$.
Here $n=1$, $f_1(w)= (w-3)^2$, $\grad f_1(w)=2(w-3)$.
Take $w_0=0$, stepsize $\alpha=0.1$ and run three iterations.

\begin{center}
\begin{tabular}{c|c|c|c}
k & $w_k$ & $v_k$ (SAGA gradient) & $f(w_k)$\\ \hline
0 & $0$ &
$\grad f_1(w_0)-\grad f_1(\phi_1^0)+\grad f_1(\phi_1^0)
 =2(0-3)-2(\phi_1^0-3)+2(\phi_1^0-3)=2(0-3)=-6$ &
$9$\\
\hline
1 & $w_1=0-0.1(-6)=0.6$ &
$\grad f_1(w_1)-\grad f_1(\phi_1^1)+\grad f_1(\phi_1^1)
 =2(0.6-3)-2(0-3)+2(0-3)=2(-2.4)-2(-3)+2(-3)= -4.8+6-6=-4.8$ &
$(0.6-3)^2=5.76$\\
\hline
2 & $w_2=0.6-0.1(-4.8)=1.08$ &
$\grad f_1(w_2)-\grad f_1(\phi_1^2)+\grad f_1(\phi_1^2)
 =2(1.08-3)-2(0.6-3)+2(0.6-3)=2(-1.92)-2(-2.4)+2(-2.4)= -3.84+4.8-4.8=-3.84$ &
$(1.08-3)^2=3.6864$\\
\hline
3 & $w_3=1.08-0.1(-3.84)=1.464$ &
\multicolumn{2}{c}{(stop after 3 iterations)}\\
\end{tabular}
\end{center}
The objective values $9\rightarrow5.76\rightarrow3.6864$ illustrate the
monotone decrease guaranteed by the theory.

\newpage
\section*{Assumptions}
\begin{assumption}[Finite sum structure]\label{ass:finite}
The objective is $f(x)=\frac1n\sum_{i=1}^{n}f_i(x)$ with each $f_i:\R^d\to\R$.
\end{assumption}
\begin{assumption}[Convexity]\label{ass:convex}
Each component $f_i$ is convex:
\[
f_i(y)\;\ge\;f_i(x)+\langle \grad f_i(x),y-x\rangle,\qquad\forall x,y\in\R^d .
\]
Consequently $f$ is convex.
\end{assumption}
\begin{assumption}[L‑smoothness]\label{ass:smooth}
There exists $L>0$ such that for all $i$ and all $x,y$,
\[
\norm{\grad f_i(x)-\grad f_i(y)}\;\le\;L\norm{x-y}.
\]
Equivalently,
\[
f_i(y)\;\le\;f_i(x)+\langle \grad f_i(x),y-x\rangle+\frac{L}{2}\norm{y-x}^2 .
\]
\end{assumption}
\begin{assumption}[Strong convexity (optional)]\label{ass:strong}
If $f$ is $\mu$‑strongly convex ($\mu>0$),
\[
f(y)\;\ge\;f(x)+\langle \grad f(x),y-x\rangle+\frac{\mu}{2}\norm{y-x}^2 .
\]
\end{assumption}
\begin{assumption}[Stepsize]\label{ass:stepsize}
The stepsize satisfies $0<\alpha\le\frac{1}{3L}$ for the convex case;
for the strongly convex case we will require
$\alpha\le\frac{1}{3(L+\mu)}$.
\end{assumption}

\section*{Main Convergence Theorem}
\begin{theorem}[Convergence of SAGA]\label{thm:convex}
Assume \ref{ass:finite}–\ref{ass:smooth} and let $\alpha\le\frac{1}{3L}$.
Define the Lyapunov function
\[
\Phi_k \;:=\; \E\!\left[ f(x_k)-f(x^\star)
    +\frac{1}{2\alpha n}\sum_{i=1}^{n}\norm{\phi_i^k-x^\star}^{2}\right],
\]
where $x^\star\in\arg\min f$. Then
\[
\Phi_{k+1}\;\le\;\left(1-\frac{\alpha\mu}{2}\right)\Phi_{k},
\qquad\text{for all }k\ge0.
\]
Consequences:
\begin{enumerate}
\item If $f$ is merely convex ($\mu=0$) we obtain the sub‑linear rate
\[
\E\!\bigl[f(\bar x_T)-f(x^\star)\bigr]\;\le\;\frac{2\bigl[f(x_0)-f(x^\star)\bigr]
    +\frac{1}{\alpha n}\sum_{i=1}^{n}\norm{\phi_i^0-x^\star}^{2}}{T},
\]
where $\bar x_T:=\frac1T\sum_{k=0}^{T-1}x_k$.
\item If $f$ is $\mu$‑strongly convex, the linear rate
\[
\E\!\bigl[f(x_k)-f(x^\star)\bigr]\;\le\;
\left(1-\frac{\alpha\mu}{2}\right)^{k}
\Bigl[ f(x_0)-f(x^\star)+\tfrac{1}{2\alpha n}\sum_{i=1}^{n}\norm{\phi_i^0-x^\star}^{2}\Bigr].
\]
\end{enumerate}
\end{theorem}

\section*{Theoretical Properties}
\begin{itemize}
\item \textbf{Variance reduction.} The term
$\grad f_{i_k}(x_k)-\grad f_{i_k}(\phi_{i_k}^k)+\bar g^k$ is an unbiased estimator
of $\grad f(x_k)$ with variance that vanishes as the table $\{\phi_i^k\}$
converges to $x^\star$.
\item \textbf{Stability w.r.t. stepsize.}
The bound $\alpha\le 1/(3L)$ ensures that the contraction derived in the proof
holds; larger stepsizes break the key inequality \eqref{ineq:lyapunov}.
\item \textbf{Comparison.}
Compared to plain SGD, SAGA enjoys an $O(1/k)$ rate for convex problems (instead of
$O(1/\sqrt{k})$) and linear convergence under strong convexity,
matching the rates of full‑gradient methods while using only a single component
gradient per iteration.
\end{itemize}

\section*{Discussion}
The theorem shows that the memory table introduces a control variate that
eliminates the stochastic noise asymptotically. The proof below follows the
classical SAGA analysis but is written with complete algebraic detail and
explicit step numbers for easy verification.

\newpage
\section*{Preliminaries (Definitions and Lemmas)}
\begin{lemma}[Smoothness inequality]\label{lem:smooth}
For any $L$‑smooth $f_i$ and any $x,y$,
\[
f_i(y) \;\le\; f_i(x) +\langle \grad f_i(x),y-x\rangle
          +\frac{L}{2}\norm{y-x}^{2}.
\]
\end{lemma}
\begin{lemma}[Convexity inequality]\label{lem:convex}
For any convex $f_i$,
\[
f_i(y) \;\ge\; f_i(x) +\langle \grad f_i(x),y-x\rangle .
\]
\end{lemma}
\begin{lemma}[Unbiasedness of the SAGA estimator]\label{lem:unbiased}
Conditioned on the sigma‑algebra $\mathcal{F}_k$ generated by
$\{x_0,\phi_i^0,\dots,x_k,\phi_i^k\}$,
\[
\E\!\bigl[ v_k \mid \mathcal{F}_k \bigr] \;=\; \grad f(x_k).
\]
\end{lemma}
\begin{proof}
By definition,
\[
v_k = \grad f_{i_k}(x_k)-\grad f_{i_k}(\phi_{i_k}^k)+\bar g^k .
\]
Taking expectation over the uniform draw of $i_k$,
\[
\E[ v_k\mid\mathcal{F}_k]
 = \frac1n\sum_{i=1}^{n}\bigl[\grad f_i(x_k)-\grad f_i(\phi_i^k)\bigr]
   +\frac1n\sum_{i=1}^{n}\grad f_i(\phi_i^k)
 = \frac1n\sum_{i=1}^{n}\grad f_i(x_k) = \grad f(x_k).
\]
\end{proof}

\section*{Main Proof (Step‑by‑Step Derivation)}
We prove Theorem~\ref{thm:convex} for the convex case ($\mu=0$); the strongly
convex case follows by adding the strong‑convexity term in Lemma~\ref{lem:convex}.
All expectations are taken with respect to the randomness up to iteration $k$.

\subsection*{Step 1: One‑step progress of the iterate}
From the update \eqref{eq:update} $x_{k+1}=x_k-\alpha v_k$ we have
\[
\norm{x_{k+1}-x^\star}^{2}
 = \norm{x_k-x^\star}^{2}
   -2\alpha\langle v_k, x_k-x^\star\rangle
   +\alpha^{2}\norm{v_k}^{2}. \tag{2}
\]
[Step 1]

\subsection*{Step 2: Take expectation and use unbiasedness}
Taking $\E[\cdot\mid\mathcal{F}_k]$ on both sides of (2) and applying Lemma~\ref{lem:unbiased},
\[
\E\!\bigl[\norm{x_{k+1}-x^\star}^{2}\mid\mathcal{F}_k\bigr]
 = \norm{x_k-x^\star}^{2}
   -2\alpha\langle \grad f(x_k),x_k-x^\star\rangle
   +\alpha^{2}\E\!\bigl[\norm{v_k}^{2}\mid\mathcal{F}_k\bigr]. \tag{3}
\]
[Step 2]

\subsection*{Step 3: Relate the inner product to function values}
By convexity of $f$ (Lemma~\ref{lem:convex}),
\[
f(x_k)-f(x^\star) \;\le\; \langle \grad f(x_k),x_k-x^\star\rangle . \tag{4}
\]
[Step 3]

\subsection*{Step 4: Upper bound the second moment of $v_k$}
Write $v_k$ as a sum of three terms:
\[
v_k = \underbrace{\grad f_{i_k}(x_k)-\grad f_{i_k}(\phi_{i_k}^k)}_{A}
      +\underbrace{\bar g^k}_{B}.
\]
Hence
\[
\norm{v_k}^{2}
 = \norm{A+B}^{2}
 \le 2\norm{A}^{2}+2\norm{B}^{2}. \tag{5}
\]
[Step 4]

\subsection*{Step 5: Bound $\E[\norm{A}^{2}\mid\mathcal{F}_k]$}
Because $i_k$ is uniform,
\[
\E\!\bigl[\norm{A}^{2}\mid\mathcal{F}_k\bigr]
 = \frac1n\sum_{i=1}^{n}\norm{\grad f_i(x_k)-\grad f_i(\phi_i^k)}^{2}
 \le \frac{L^{2}}{n}\sum_{i=1}^{n}\norm{x_k-\phi_i^k}^{2}  \tag{6}
\]
where the inequality follows from $L$‑smoothness (Lemma~\ref{lem:smooth}).
[Step 5]

\subsection*{Step 6: Bound $\norm{B}^{2}$}
Recall $B=\bar g^k = \frac1n\sum_{i=1}^{n}\grad f_i(\phi_i^k)$. By Jensen’s inequality,
\[
\norm{B}^{2}
 \le \frac1n\sum_{i=1}^{n}\norm{\grad f_i(\phi_i^k)}^{2}
 \le \frac{L^{2}}{n}\sum_{i=1}^{n}\norm{\phi_i^k-x^\star}^{2}
    +\frac{2L}{n}\sum_{i=1}^{n}\bigl[f(\phi_i^k)-f(x^\star)\bigr]. \tag{7}
\]
The first term uses the $L$‑smoothness bound
$\norm{\grad f_i(\phi_i^k)}\le L\norm{\phi_i^k-x^\star}$,
the second term follows from the convexity inequality
$f_i(\phi_i^k)-f_i(x^\star)\ge\langle\grad f_i(\phi_i^k),\phi_i^k-x^\star\rangle$ and the
Cauchy–Schwarz inequality.
[Step 6]

\subsection*{Step 7: Assemble the bound on $\E[\norm{v_k}^{2}\mid\mathcal{F}_k]$}
Combining (5)–(7) gives
\[
\E\!\bigl[\norm{v_k}^{2}\mid\mathcal{F}_k\bigr]
 \le \frac{2L^{2}}{n}\sum_{i=1}^{n}\norm{x_k-\phi_i^k}^{2}
   +\frac{2L^{2}}{n}\sum_{i=1}^{n}\norm{\phi_i^k-x^\star}^{2}
   +\frac{4L}{n}\sum_{i=1}^{n}\bigl[f(\phi_i^k)-f(x^\star)\bigr]. \tag{8}
\]
[Step 7]

\subsection*{Step 8: Lyapunov function}
Define
\[
\Phi_k \;:=\; \E\!\bigl[ f(x_k)-f(x^\star)
    +\frac{1}{2\alpha n}\sum_{i=1}^{n}\norm{\phi_i^k-x^\star}^{2}\bigr]. \tag{9}
\]
Our goal is to show $\Phi_{k+1}\le\Phi_k$ up to a contraction factor.

\subsection*{Step 9: Expected decrease of the table term}
Only the entry $i_k$ changes:
\[
\phi_{i_k}^{k+1}=x_k,\qquad
\phi_j^{k+1}=\phi_j^{k}\;(j\neq i_k).
\]
Hence
\[
\E\!\Bigl[\sum_{i=1}^{n}\norm{\phi_i^{k+1}-x^\star}^{2}\mid\mathcal{F}_k\Bigr]
 = \sum_{i=1}^{n}\norm{\phi_i^{k}-x^\star}^{2}
   -\frac1n\sum_{i=1}^{n}\norm{\phi_i^{k}-x^\star}^{2}
   +\frac1n\sum_{i=1}^{n}\norm{x_k-x^\star}^{2}. \tag{10}
\]
The first subtraction accounts for the replaced entry, the last term is the
new contribution of the sampled index.
[Step 9]

\subsection*{Step 10: Combine iterate and table recursions}
Multiply (10) by $\frac{1}{2\alpha n}$ and add it to (3) after taking total
expectation. Using (4) to replace the inner product term and (8) for the
second‑moment term, we obtain
\[
\begin{aligned}
\E\!\bigl[\Phi_{k+1}\mid\mathcal{F}_k\bigr]
 &\le \Phi_k
   -2\alpha\bigl[f(x_k)-f(x^\star)\bigr]      \\
 &\quad +\alpha^{2}\Bigl[
        \frac{2L^{2}}{n}\sum_{i=1}^{n}\norm{x_k-\phi_i^k}^{2}
       +\frac{2L^{2}}{n}\sum_{i=1}^{n}\norm{\phi_i^k-x^\star}^{2}
       +\frac{4L}{n}\sum_{i=1}^{n}\bigl[f(\phi_i^k)-f(x^\star)\bigr]
      \Bigr] \\
 &\quad -\frac{1}{2\alpha n}\Bigl[
        \frac1n\sum_{i=1}^{n}\norm{\phi_i^{k}-x^\star}^{2}
        -\norm{x_k-x^\star}^{2}
      \Bigr].
\end{aligned} \tag{11}
\]
[Step 10]

\subsection*{Step 11: Bounding the mixed term $\norm{x_k-\phi_i^k}^{2}$}
By the elementary identity
$\norm{x_k-\phi_i^k}^{2}
 = \norm{x_k-x^\star}^{2}
   +\norm{\phi_i^k-x^\star}^{2}
   -2\langle x_k-x^\star,\phi_i^k-x^\star\rangle$,
and using $2\langle a,b\rangle\le \norm{a}^{2}+\norm{b}^{2}$,
\[
\norm{x_k-\phi_i^k}^{2}
 \le 2\norm{x_k-x^\star}^{2}+2\norm{\phi_i^k-x^\star}^{2}. \tag{12}
\]
[Step 11]

\subsection*{Step 12: Insert (12) into (11) and simplify}
After substituting (12) and collecting like terms,
\[
\begin{aligned}
\E\!\bigl[\Phi_{k+1}\mid\mathcal{F}_k\bigr]
 &\le \Phi_k
   -2\alpha\bigl[f(x_k)-f(x^\star)\bigr] \\
 &\quad +\alpha^{2}\frac{4L^{2}}{n}\sum_{i=1}^{n}\norm{x_k-x^\star}^{2}
   +\alpha^{2}\frac{4L^{2}}{n}\sum_{i=1}^{n}\norm{\phi_i^k-x^\star}^{2}\\
 &\quad +\alpha^{2}\frac{4L}{n}\sum_{i=1}^{n}\bigl[f(\phi_i^k)-f(x^\star)\bigr] \\
 &\quad -\frac{1}{2\alpha n^{2}}\sum_{i=1}^{n}\norm{\phi_i^{k}-x^\star}^{2}
   +\frac{1}{2\alpha n}\norm{x_k-x^\star}^{2}.
\end{aligned} \tag{13}
\]
[Step 12]

\subsection*{Step 13: Choose stepsize to make the coefficient of
$\norm{x_k-x^\star}^{2}$ non‑positive}
Recall that for $L$‑smooth convex $f$,
\[
f(x_k)-f(x^\star)\;\ge\;\frac{1}{2L}\norm{\grad f(x_k)}^{2}
\;\ge\;0,
\]
so we may drop the non‑negative term $-2\alpha[f(x_k)-f(x^\star)]$ if needed.
To guarantee a contraction we enforce
\[
-2\alpha +\alpha^{2}\frac{4L^{2}}{n} +\frac{1}{2\alpha n}\;\le\;0.
\]
A sufficient condition is $\alpha\le\frac{1}{3L}$ (Assumption~\ref{ass:stepsize}).
Under this choice the coefficient of $\norm{x_k-x^\star}^{2}$ becomes
non‑positive, and (13) reduces to
\[
\E\!\bigl[\Phi_{k+1}\mid\mathcal{F}_k\bigr]
 \le \Phi_k
    -\underbrace{\Bigl(\alpha -\frac{2\alpha^{2}L}{n}\Bigr)}_{\ge 0}
      \frac{1}{n}\sum_{i=1}^{n}\bigl[f(\phi_i^k)-f(x^\star)\bigr]. \tag{14}
\]
[Step 13]

\subsection*{Step 14: Take full expectation}
Taking expectation over $\mathcal{F}_k$ in (14) yields
\[
\E[\Phi_{k+1}] \;\le\; \E[\Phi_k] . \tag{15}
\]
Iterating (15) gives $\E[\Phi_T]\le\Phi_0$ for any $T\ge1$.
Expanding $\Phi_T$ and using the definition (9) we obtain
\[
\E\!\bigl[f(x_T)-f(x^\star)\bigr]
 \le \frac{\Phi_0}{T}
 = \frac{ f(x_0)-f(x^\star)
        +\frac{1}{2\alpha n}\sum_{i=1}^{n}\norm{\phi_i^0-x^\star}^{2}}
        {T}. \tag{16}
\]
[Step 14]

\subsection*{Step 15: Strongly convex case}
If $f$ is $\mu$‑strongly convex, Lemma~\ref{lem:convex} improves (4) to
\[
\langle \grad f(x_k),x_k-x^\star\rangle
 \;\ge\; f(x_k)-f(x^\star)+\frac{\mu}{2}\norm{x_k-x^\star}^{2}. \tag{17}
\]
Repeating Steps 1–13 with (17) gives an extra term $-\alpha\mu\norm{x_k-x^\star}^{2}$
that survives after the stepsize restriction. Consequently,
\[
\E[\Phi_{k+1}] \;\le\;
\Bigl(1-\frac{\alpha\mu}{2}\Bigr)\E[\Phi_k],
\]
which after unwinding yields the linear rate claimed in Theorem~\ref{thm:convex}.
[Step 15]

\section*{Rate Derivation (With Constants)}
\begin{itemize}
\item \textbf{Convex case.} From (16) and the definition of $\Phi_0$,
\[
\E\!\bigl[f(\bar x_T)-f(x^\star)\bigr]
 \le \frac{2\bigl[f(x_0)-f(x^\star)\bigr]
          +\frac{1}{\alpha n}\sum_{i=1}^{n}\norm{\phi_i^0-x^\star}^{2}}{T}.
\]
Choosing $\alpha=1/(3L)$ gives an explicit constant $6L$ in front of the
initial error.
\item \textbf{Strongly convex case.} With $\alpha\le 1/(3(L+\mu))$,
\[
\E\!\bigl[f(x_k)-f(x^\star)\bigr]
 \le \Bigl(1-\frac{\alpha\mu}{2}\Bigr)^{k}
    \Bigl[ f(x_0)-f(x^\star)+\frac{1}{2\alpha n}\sum_{i=1}^{n}\norm{\phi_i^0-x^\star}^{2}\Bigr].
\]
Setting $\alpha=1/(3(L+\mu))$ yields contraction factor
$1-\frac{\mu}{6(L+\mu)}$.
\end{itemize}

\section*{Special Cases}
\begin{enumerate}
\item \textbf{Full‑gradient regime ($n=1$).} The table contains a single entry,
so $v_k=\grad f_1(x_k)$ and SAGA reduces to vanilla gradient descent with
stepsize $\alpha\le 1/(3L)$, reproducing the classical $O(1/k)$ and linear
rates.
\item \textbf{Large‑scale regime ($n\gg 1$).} The variance reduction term
$\frac{1}{n}\sum_i\grad f_i(\phi_i^k)$ approximates the full gradient after a
few epochs, thus the per‑iteration cost stays $\mathcal{O}(d)$ while the
convergence matches that of the full gradient method.
\item \textbf{Non‑strongly convex but smooth.} The $O(1/k)$ bound holds without
any additional regularization; adding a small $\ell_2$ term makes the problem
strongly convex and immediately yields linear convergence.
\end{enumerate}

\end{document}